\section{Heat flow and the complementary physical space}\label{sec:heat}

The second manuscript of the latest edition, \texttt{HEAT\_AND\_COMPLEMENT.md} (H1--H34; Appendix~B of the fifth-reference reader), uses the control of the vacuum from Section~\ref{sec:gap} for two further purposes: box-independent bounds on the heat correlations of plaquettes, and bounds on how much the plaquette family mixes with the rest of the physical space. This section states both, with proofs of the steps that carry the constants.

\subsection{Objects}

Let $A=H_L-E_0$ on the centred physical space and $\rho=\psi^2$. Every statement below holds for every $L\ge2$ and $a>0$.
\begin{itemize}
\item \emph{Centred plaquette states and their heat correlations}:
\[
r_p=(W_p-\langle W_p\rangle_\rho)\psi,\qquad \widehat C_{pq}(\tau;\xi)=\langle r_p,e^{-\tau A/\kappa}r_q\rangle,\qquad\tau=\kappa t.
\]
\item \emph{Row norm}: $\|B\|_{\mathrm{row}}=\max_p\sum_q|B_{pq}|$.
\item \emph{Taylor coefficients}: $C_0,C_2,C_4$ are the coefficients of $\xi^0,\xi^2,\xi^4$ in $\widehat C$; the workbench computed them in its 21 September heat edition.
\item \emph{The transformed generator}: in the coordinates $f\mapsto\psi f$, Lemma~\ref{lem:gst} turns $A/\kappa$ into $K-D_v$ with $D_vh=2\Gamma(v,h)$.
\item \emph{The Fourier-algebra norm} is $\|\cdot\|_X$ of Section~\ref{sec:gap}, and $\chi=4\sum_it_i\xi^i+\frac23w_*$ is the number from the proof of Proposition~\ref{prop:return}, so $3(1-\chi)=d_5(\xi)$ and $\|D_vh\|_X\le\chi\|Kh\|_X$.
\end{itemize}

\subsection{A heat remainder independent of the box}

\begin{theorem}[heat remainder; H13]\label{thm:heat}
For every $L\ge2$, every $\tau\ge0$ and every $|\xi|<1/55$,
\[
\bigl\|\widehat C(\tau;\xi)-C_0(\tau)-\xi^2C_2(\tau)-\xi^4C_4(\tau)\bigr\|_{\mathrm{row}}\le\frac{2829}{13}\,e^{-13\tau/8}\,\frac{(55|\xi|)^6}{1-(55|\xi|)^2}.
\]
The workbench states this with the prefactor $67896/169=401.75\ldots$ in place of $2829/13=217.61\ldots$; its bound is also valid.
\end{theorem}
\status{\textsf{Conditional}; \textsf{audited in part}.
\begin{itemize}
\item \emph{Conditional} on the analytic core of Theorem~\ref{thm:ym01}, and therefore on the same order-3 to order-5 inputs.
\item \emph{Checked here.} Every constant is re-derived below. The evenness and the Cauchy step are proved.
\item \emph{Read, not checked in every step.} The semigroup construction (H5--H7) and the analytic continuation to complex $\xi$ are standard. Proposition~\ref{prop:complexgap} gives the spectral side of that continuation.
\item \emph{The prefactor.} It is sharpened here by the mean estimate of step (i), which was found by the referee pass of version~1.
\item \emph{Not needed.} The inherited matrices $C_0,C_2,C_4$ are not needed for the truth of the bound, only for its use.
\item \emph{The earlier edition's bounds.}
\begin{itemize}
\item Its bound on the circle $3/256$, with prefactor $512$ and rate $3/2$, is superseded.
\item Its second bound U29b (heat reader of 21 September, p.~37), on the circle $45/4096$ with prefactor $6184/25$ and rate $15/8$, is sharper than the present theorem for small $|\xi|$ once $\tau>12.6$ ($\tau>10.15$ with the old prefactor). It is superseded by the family of radii below, which at $R=45/4096$ has rate $2.2588$ and prefactor $85.0$.
\end{itemize}
\end{itemize}}

The proof is in four steps; the constants are collected in Proposition~\ref{prop:heatconst}.

\emph{(i) The mean and the semigroup} (H1--H7). Put $T=Q_HD_vK^{-1}$ and $p=P_HD_vK^{-1}$ on the nonconstant physical coefficients. Since the norm $\|\cdot\|_X$ adds the scalar and the nonconstant parts, $|py|+\|Ty\|_X\le\chi\|y\|_X$. So $S=(I-T)^{-1}$ exists with $\|S\|\le1/(1-\chi)$, and the mean is $\mu(F)=P_HF+pSQ_HF$. The workbench bounds $|(\mu-P_H)F|$ by $\chi\|Q_HF\|_X/(1-\chi)$. The additive form of the inequality gives more:
\[
|(\mu-P_H)F|\le\chi\|Q_HF\|_X,\qquad\text{so}\qquad|\mu F|\le\|F\|_X\ \text{for every }\chi<1.
\]
Indeed, let $y=Q_HF$ and $z=Sy$, so that $z=y+Tz$. Then
\[
|pz|\le\chi\|z\|_X-\|Tz\|_X\le\chi(\|y\|_X+\|Tz\|_X)-\|Tz\|_X=\chi\|y\|_X-(1-\chi)\|Tz\|_X\le\chi\|y\|_X .
\]
(Found by the referee pass of version~1; checked here, also on random examples in $\ell^1$.)
For $L_0=(I-T)K$ and $h>0$, coefficientwise $\|(I+hK)f\|_X=\|f\|_X+h\|Kf\|_X$. Hence
\[
\|(I+hL_0)f\|_X\ge\|f\|_X+h(1-\chi)\|Kf\|_X\ge(1+3h(1-\chi))\|f\|_X,
\]
using $c_j\ge3$ on nonconstant physical blocks (Lemma~\ref{lem:casimir}). With the range condition (H5), the Lumer--Phillips theorem \cite{LumerPhillips} gives a semigroup $E(\tau)=e^{-\tau L_0}$ with $\|E(\tau)\|\le e^{-3(1-\chi)\tau}$.

\emph{(ii) A formula for the correlations.} Let $\mathcal T(\tau)$ be the transformed heat semigroup, so that $\widehat C_{pq}(\tau)=\mu[(\mathcal T(\tau)F_p)W_q]$ with $F_p=W_p-\mu W_p$; here $\mu[\mathcal T(\tau)F_p]=\mu[F_p]=0$ removes the second mean. Let $h_\tau=K^{-1}SE(\tau)W_p$ and $G_z=\sum_qz_qW_q$ with $|z_q|\le1$. Since $W_p$ has Haar mean $0$, the Poisson equation (H3) gives $(K-D_v)h_\tau=E(\tau)W_p-\mu(E(\tau)W_p)=\mathcal T(\tau)F_p$. For any $h$ and $G$,
\[
\int e^{2v}\bigl(Kh-2\Gamma(v,h)\bigr)G\,dU=\int e^{2v}\,\Gamma(h,G)\,dU,
\]
by moving one derivative $X_{e,a}$ onto $e^{2v}G$. Dividing by $\int e^{2v}$ gives $\sum_qz_q\widehat C_{pq}(\tau)=\mu\Gamma(h_\tau,G_z)$, and $\|Kh_\tau\|_X\le\|S\|\,\|E(\tau)\|\,\|W_p\|_X\le\frac8{1-\chi}e^{-3(1-\chi)\tau}$.

\begin{proposition}[the constants of the heat bound; H9--H12]\label{prop:heatconst}
In this setting:
\begin{enumerate}[label=(\alph*)]
\item $\|\Gamma(h,G_z)\|_X\le32\|Kh\|_X$;
\item $|P_H\Gamma(h,G_z)|\le\frac18\|Kh\|_X$;
\item $\|\widehat C(\tau)\|_{\mathrm{row}}\le\dfrac{1+255\chi}{1-\chi}\,e^{-3(1-\chi)\tau}$, and a fortiori the workbench's $\dfrac{1+255\chi}{(1-\chi)^2}\,e^{-3(1-\chi)\tau}$;
\item at $|\xi|=1/55$ one has $1/55<\alpha_5$ and $d_5(1/55)=1.637\ldots>13/8$, so $\chi<11/24$. The prefactor is at most $\frac{1+255\cdot11/24}{13/24}=\frac{2829}{13}$ (the workbench's $\frac{67896}{169}$ with its form of (c)), and the rate is at least $\frac{13}8$.
\end{enumerate}
\end{proposition}
\status{Proved here. The workbench's (H9)--(H12) state (a), (b), the weaker form of (c) and (d). The earlier heat reader of 21 September already derives $32$ (its U28), the rule $\frac18$, and the weaker form of (c) (its U29a, pp.~36--37); this paper writes out the derivation. The sharper form of (c) uses the mean estimate of step (i). The inequalities of (d) are checked in exact rational arithmetic (scripts \note{ymr\_heat\_constants.py} and \note{ymr\_referee\_checks.py}).}
\begin{proof}
(a) By (F10), $\|\Gamma(h,G)\|_X\le2t(G)\|Kh\|_X$. The function $W_q$ has one Fourier block, with spin $\frac12$ on the four links of $q$, and its coefficient matrix has trace norm $2^{4-1}=8$ by (F12): four fundamental occurrences, two cyclic sign changes. A link lies on at most four plaquettes, so $t(G_z)\le4\cdot\frac12\cdot8=16$.
(b) Integration by parts gives $P_H\Gamma(h,G_z)=\int h\,KG_z\,dU=3\sum_qz_q\int hW_q\,dU$, since each $W_q$ has Casimir $4\cdot\frac34=3$. The gauge-invariant functions whose Fourier support is spin $\frac12$ on the four links of $q$ and spin $0$ elsewhere form the line $\C W_q$: at each corner two spin-$\frac12$ links meet, and $\frac12\otimes\frac12$ contains the trivial representation once. Write $a_q$ for the component of $h$ on this line. By Haar orthogonality and $\int W_q^2=1$, $\int hW_q=a_q$. The block of $a_qW_q$ contributes $3\cdot8|a_q|$ to $\|Kh\|_X$, so $|P_H\Gamma|\le3\sum|a_q|\le\frac18\|Kh\|_X$.
(c) $\mu\Gamma=P_H\Gamma+(\mu-P_H)Q_H\Gamma$, and $\|Q_H\Gamma\|_X=\|\Gamma\|_X-|P_H\Gamma|$. By step (i), then (b) and (a),
\begin{align*}
|\mu\Gamma(h_\tau,G_z)|&\le(1-\chi)|P_H\Gamma|+\chi\|\Gamma\|_X\le\Bigl(\frac{1-\chi}8+32\chi\Bigr)\|Kh_\tau\|_X\\
&\le\Bigl(\frac{1-\chi}8+32\chi\Bigr)\frac{8e^{-3(1-\chi)\tau}}{1-\chi}=\frac{1+255\chi}{1-\chi}e^{-3(1-\chi)\tau}.
\end{align*}
With the workbench's form of step (i), the same steps give $\bigl(\frac18+\frac{32\chi}{1-\chi}\bigr)\frac8{1-\chi}=\frac{1+255\chi}{(1-\chi)^2}$. The row norm is the maximum over the phases $z_q$ of the left side.
(d) The value $d_5(1/55)$ is computed from the rationals of (F26), and $d_5>13/8$ is equivalent to $\chi<11/24$ since $d_5=3(1-\chi)$. Both prefactors increase with $\chi$, and $3(1-\chi)>3\cdot\frac{13}{24}=\frac{13}8$.
\end{proof}

\emph{(iii) Evenness.} The map $U_{(n,i)}\mapsto(-1)^{\sum_{k<i}n_k}U_{(n,i)}$ multiplies each link by a central element. It preserves the Haar measure, $K$ and the gauge action, and it sends every $W_p$ to $-W_p$: for $p=(n;i,j)$ with $i<j$ the four signs multiply to $-1$. It therefore conjugates $H_L(\xi)$ to $H_L(-\xi)$ plus a constant, and maps each $r_p$ to $-r_p$ at the opposite coupling. So $\widehat C(\tau;-\xi)=\widehat C(\tau;\xi)$, and only even powers of $\xi$ occur. (Checked on all $240$ faces of the box $L=2$.)

\emph{(iv) Cauchy's estimate.} By (c) and (d), on the circle $|\zeta|=1/55$ every phase-weighted row sum of $\widehat C(\tau;\zeta)$ is at most $Ce^{-13\tau/8}$ with $C=2829/13$. So the coefficient of $\zeta^{2n}$ has row norm at most $C e^{-13\tau/8}55^{2n}$, and summing over $n\ge3$ gives the stated tail $\sum_{n\ge3}(55|\xi|)^{2n}=(55|\xi|)^6/(1-(55|\xi|)^2)$. This proves Theorem~\ref{thm:heat}.

The same step gives, for every $n$, $\|[\xi^{2n}]\widehat C(\tau)\|_{\mathrm{row}}\le\frac{2829}{13}\,55^{2n}e^{-13\tau/8}$, and the analogous remainder after any even order. At $\xi=0$ one has $\psi=1$ and $r_p=W_p$, an eigenfunction of $K$ with eigenvalue $3$, so $C_0(\tau)=e^{-3\tau}I$ exactly. (Both noted by the referee pass.)

\emph{Other radii} (\textsf{Generalised here}). Nothing in steps (i)--(iv) uses $R=1/55$ except the numbers of Proposition~\ref{prop:heatconst}(d). For every $0<R\le\alpha_5$ the same proof gives the bound of Theorem~\ref{thm:heat} with the following replacements: $55|\xi|$ by $|\xi|/R$; the prefactor by $C(R)=(1+255\chi_R)/(1-\chi_R)$; the rate $13/8$ by $d_5(R)=3(1-\chi_R)$, where $\chi_R=1-d_5(R)/3$. At the endpoint $R=\alpha_5$, apply the Cauchy step at radii $r<R$ and let $r$ increase to $R$, as the earlier edition does.
\begin{itemize}
\item At $R=1/55$ the exact $\chi_R=0.45412\ldots$ gives the prefactor $213.97\ldots$ and the rate $1.6376\ldots$; the rounding $\chi<11/24$ costs little.
\item At $R=\alpha_5$ the prefactor is $241.99\ldots$ and the rate $1.5453\ldots$.
\item At $R=45/4096$, the radius of U29b, the prefactor is $85.00\ldots$ and the rate $2.2588\ldots$.
\end{itemize}
With the workbench's form $(1+255\chi_R)/(1-\chi_R)^2$ these prefactors are $391.98$, $469.79$ and $112.89$. They are computed from the rationals of (F26); they are floating-point values, not rational certificates.

\subsection{Relaxation into the whole complementary space}\label{ssec:complement}

Let $R$ be the map with columns $r_p$ and $\Phi=\kappa A^{-1}R$. The original Grams and the energy of the family are
\[
G_0=R^*R,\quad G_1=R^*\Phi,\quad G_2=\Phi^*\Phi,\quad E=\kappa G_1,\quad K_o=\kappa^{-1}R^*AR.
\]
Put $P=\Phi G_2^{-1}\Phi^*$, $Q=I-P$, $B_\Phi=QR$ and $D_Q=QAQ$ on $Q\mathcal H_0$. Allowing the family $\Phi c$ to relax into every $h\in Q\mathcal H_0$ gives the relaxed energy and Gram
\[
E_{\mathrm{full}}=E-\kappa^2B_\Phi^*D_Q^{-1}B_\Phi,\qquad G_{\mathrm{full}}=G_2+\kappa^2B_\Phi^*D_Q^{-2}B_\Phi .
\]

\begin{proposition}[relaxation bounds; H20--H29]\label{prop:relax}
Suppose that, for numbers $w_0,w_1,w_2,w_K,\beta\ge0$ and $d_{\mathrm{ph}}>0$,
\[
\|G_k-3^{-k}I\|\le w_k\ (k=0,1,2),\qquad\|K_o-3I\|\le w_K,\qquad 0\preceq J=G_0-6G_1+9G_2\preceq\beta I,\qquad A\ge\kappa d_{\mathrm{ph}}.
\]
Put $l_k=3^{-k}-w_k>0$, $u_k=3^{-k}+w_k$ and $u_K=3+w_K$. Then:
\begin{enumerate}[label=(\alph*)]
\item the fraction of each $\Phi c$ seen by the plaquette states, $\|P_R\Phi c\|^2/\|\Phi c\|^2$, is at least $\max\{l_1^2/(u_0u_2),\,1-\beta/(9l_2)\}$, where $P_R$ is the orthogonal projection onto the span of the $r_p$;
\item $(1-\eta_E)E\preceq E_{\mathrm{full}}\preceq E$ and $G_2\preceq G_{\mathrm{full}}\preceq(1+\eta_G)G_2$, with $\eta_E=\beta/(d_{\mathrm{ph}}l_1)$ and $\eta_G=\beta/(d_{\mathrm{ph}}^2l_2)$;
\item the energy of $(I-P_R)\Phi c$ is at most $(u_1u_K/l_0^2-1)$ times $c^*Ec$.
\end{enumerate}
\end{proposition}
\status{Proved here; the workbench's (H21), (H23), (H26)--(H29), with the second bound of (a) added. The domain statement for $D_Q$ (H25) is the workbench's and was read.}
\begin{proof}
(a) $\|P_R\Phi c\|^2=c^*G_1G_0^{-1}G_1c\ge u_0^{-1}c^*G_1^2c\ge(l_1^2/u_0)|c|^2$ and $\|\Phi c\|^2=c^*G_2c\le u_2|c|^2$. For the second bound (found by the referee pass): $G_1=\kappa R^*A^{-1}R$ is symmetric, so $J=(R-3\Phi)^*(R-3\Phi)$, and $\|(I-P_R)\Phi c\|^2\le\|\Phi c-\frac13Rc\|^2=\frac19c^*Jc\le\frac19\beta|c|^2$, while $\|\Phi c\|^2\ge l_2|c|^2$.
(b) $B_\Phi^*B_\Phi=G_0-G_1G_2^{-1}G_1$ is the least value of $(R-\Phi X)^*(R-\Phi X)$ over coefficient matrices $X$, and $X=3I$ gives $J$; so $B_\Phi^*B_\Phi\preceq J\preceq\beta I$. Since $D_Q\succeq\kappa d_{\mathrm{ph}}$, $\kappa^2B_\Phi^*D_Q^{-1}B_\Phi\preceq(\kappa/d_{\mathrm{ph}})\beta I\preceq(\beta/(d_{\mathrm{ph}}l_1))\kappa G_1$, and $\kappa^2B_\Phi^*D_Q^{-2}B_\Phi\preceq(\beta/d_{\mathrm{ph}}^2)I\preceq(\beta/(d_{\mathrm{ph}}^2l_2))G_2$.
(c) Since $A\Phi=\kappa R$ and $P_RR=R$, the mixed energy is $q_A(\Phi c,P_R\Phi c)=\kappa\langle Rc,\Phi c\rangle=\kappa c^*G_1c=q_A(\Phi c)$, and $P_R\Phi c=RG_0^{-1}G_1c$ has energy $\kappa c^*G_1G_0^{-1}K_oG_0^{-1}G_1c$. Expanding $q_A((I-P_R)\Phi c)$ gives $\kappa c^*(G_1G_0^{-1}K_oG_0^{-1}G_1-G_1)c$, the workbench's (H29). Now $G_1G_0^{-1}K_oG_0^{-1}G_1\preceq(u_K/l_0^2)G_1^2\preceq(u_Ku_1/l_0^2)G_1$.
\end{proof}

The workbench obtains the numbers $w_k$, $w_K$ and $\beta$ from Theorem~\ref{thm:heat}, and $w_K$ from a bound on the kinetic rows, which carries the constants $30$ and $48$ (H14).

\begin{proposition}[the kinetic rows; H14]\label{prop:kinetic}
For each plaquette $p$ and each complex $|\zeta|\le1/55$, the entries $\mu\Gamma(W_p,W_q)$ of the kinetic matrix satisfy:
\begin{enumerate}[label=(\alph*)]
\item $\Gamma(W_p,W_p)=3-\chi_1(U_p)$, of norm $\|\cdot\|_X=3+27=30$;
\item for each of the $12$ plaquettes $q$ adjacent to $p$, $\Gamma(W_p,W_q)=\frac34P_0-\frac14P_1$. Here $P_0,P_1$ are the parts of $W_pW_q$ with spin $0$ and spin $1$ on the shared link, of Casimirs $\frac92$ and $\frac{13}2$. So $\|\Gamma(W_p,W_q)\|_X$ is $24$ for $8$ of the $q$ and $6+6\sqrt3$ for the other $4$; both are below $48$. For non-adjacent $q\ne p$, $\Gamma(W_p,W_q)=0$;
\item hence $\sum_q|\mu\Gamma(W_p,W_q)|\le3+\chi(27+192+24+24\sqrt3)<134$ for $\chi<\frac{11}{24}$. The workbench's bound is $30+12\cdot48=606$.
\end{enumerate}
\end{proposition}
\begin{proof}
From $2\Gamma(f,h)=(c_f+c_h)fh-K(fh)$ for block eigenfunctions (the workbench's F17):
\begin{itemize}
\item \emph{Diagonal.} $W_p^2=1+\chi_1(U_p)$ and $\chi_1(U_p)$ has Casimir $4\cdot2=8$, so $2\Gamma(W_p,W_p)=6(1+\chi_1)-8\chi_1$. The norm of $\chi_1(U_p)$ is $3^3=27$ by F12 with $d=3$.
\item \emph{Adjacent faces.} $2\Gamma(W_p,W_q)=6(P_0+P_1)-\frac92P_0-\frac{13}2P_1$.
\item \emph{The block norms.} The trace norms of $P_0$ and $P_1$ were computed from explicit $SU(2)$ tensors in the first-pass audit (\note{46a\_}). They are $(16,48)$ for two thirds of the adjacent pairs and $(8,24\sqrt3)$ for the rest, which gives $\frac34\cdot16+\frac14\cdot48=24$ and $\frac34\cdot8+\frac14\cdot24\sqrt3=6+6\sqrt3$. The referee pass of this paper found the split $8+4$ for the twelve neighbours of one face.
\item \emph{Non-adjacent faces.} For non-adjacent $q\ne p$ no link is shared, so no derivative acts on both factors, and $\Gamma=0$.
\item \emph{The row.} Only the diagonal entry has a constant part, namely $3$. So step (i) gives $|\mu\Gamma|\le|P_H\Gamma|+\chi\|Q_H\Gamma\|_X$ entrywise, and the sum follows.
\end{itemize}
\end{proof}
\status{Proved here, given the first pass's block norms; the sharper row sum found by the referee pass. The workbench's constants $30$ and $48$ are thereby derived, and $48$ is not attained.} The time-integral formula $G_k=\int_0^\infty\tau^{k-1}\widehat C(\tau)\,d\tau/(k-1)!$ ($k\ge1$) gives the remainder factors $C/d^k$ with $d=13/8$. For $J=\widehat C(0)+\int_0^\infty(9\tau-6)\widehat C(\tau)\,d\tau$ the signed weight is kept: $\int_0^\infty|9\tau-6|e^{-d\tau}d\tau=\frac6d-\frac9{d^2}+\frac{18}{d^2}e^{-2d/3}$, and with $e^{-13/12}<\frac{339}{1000}$ the prefactor is $C_J=5156090136/3570125$ (H17--H18). The degree-2 and degree-4 parts are the workbench's inherited row bounds (the table after H18), which come from 199 supports and 559 marked contributions of its earlier volume edition.

\begin{proposition}[the benchmark table; H5]\label{prop:bench}
From the stated inherited row bounds and the constants $C$, $C_J$ and $606$, the widths and Proposition~\ref{prop:relax} give, for every $L\ge2$ and $a>0$:

\begin{center}\small
\begin{tabular}{@{}lccc@{}}
\toprule
Coupling & Observed fraction & Energy loss $\eta_E$ & Gram increase $\eta_G$\\
\midrule
$g^2\ge10$ & $>0.9778$ & $<0.01057$ & $<0.01123$\\
$g^2\ge12$ & $>0.9975$ & $<0.00115$ & $<0.00120$\\
$g^2\ge25/2$ & $>0.9984$ & $<0.00071$ & $<0.00073$\\
$g^2\ge13$ & $>0.99904$ & $<0.000436$ & $<0.000451$\\
$g^2\ge16$ & $>0.99991$ & $<0.0000357$ & $<0.0000364$\\
\bottomrule
\end{tabular}
\end{center}
In particular the workbench's table holds: at $g^2\ge13$ the energy loss and the Gram increase are below $1/2000$, and at $g^2\ge16$ below $1/25000$.
\end{proposition}
\status{\textsf{Reproduced}: recomputed in exact rational arithmetic from the stated inputs (script \note{ymr\_heat\_constants.py}); at $g^2=13$, with the workbench's $d_{\mathrm{ph}}=29047/10000$, the three values agree with its diagnostics to ten digits. At the other couplings the script uses $d_{\mathrm{ph}}$ equal to $d_5(\xi)$ rounded down to four digits. \textsf{Conditional} on the inherited degree-2 and degree-4 row bounds and on the inputs of Theorem~\ref{thm:ym01}. The widths are monotone in $\xi$ and $d_5$ decreases in $\xi$, so each row holds for all larger $g^2$, as the workbench says.}

With the sharper constants of this paper the same computation gives more. These are the prefactor $2829/13$ of Theorem~\ref{thm:heat}, the leakage prefactor $C_J=782.29\ldots$ (the workbench's $C_J$ times $13/24$), the kinetic bound $134$ of Proposition~\ref{prop:kinetic}, and the second bound of Proposition~\ref{prop:relax}(a):

\begin{center}\small
\begin{tabular}{@{}lccc@{}}
\toprule
Coupling & Observed fraction & Energy loss $\eta_E$ & Gram increase $\eta_G$\\
\midrule
$g^2\ge10$ & $>0.99458$ & $<0.005713$ & $<0.006053$\\
$g^2\ge12$ & $>0.99940$ & $<0.000623$ & $<0.000647$\\
$g^2\ge25/2$ & $>0.99963$ & $<0.000380$ & $<0.000394$\\
$g^2\ge13$ & $>0.99977$ & $<0.000237$ & $<0.000244$\\
$g^2\ge16$ & $>0.999981$ & $<0.0000194$ & $<0.0000198$\\
\bottomrule
\end{tabular}
\end{center}
(Script \note{ymr\_referee\_checks.py}; the referee pass computed the same values independently.)

The plaquette states also bound the gap from above. This was found by the referee pass; the proof was checked here.

\begin{proposition}[an upper bound on the gap]\label{prop:gapupper}
\begin{enumerate}[label=(\alph*)]
\item For every $L\ge2$, $a,g>0$ and every plaquette $p$,
\[
\Delta_L\le\kappa\,\frac{3-\langle\chi_1(U_p)\rangle_\rho}{1+\langle\chi_1(U_p)\rangle_\rho-\langle W_p\rangle_\rho^2}.
\]
\item With the widths of Proposition~\ref{prop:relax}, $\Delta_L/\kappa\le(3+w_K)/(1-w_0)$. With the constants above, the gap therefore lies in
\begin{itemize}
\item $[2.8382,\,3.0055]\kappa$ for $g^2\ge10$;
\item $[2.9047,\,3.00024]\kappa$ for $g^2\ge13$;
\item $[2.9372,\,3.00003]\kappa$ for $g^2\ge16$.
\end{itemize}
The lower ends are $d_5(\xi)$, under the inputs of Theorem~\ref{thm:ym01}.
\item At each fixed $L$, the first-order term of $\Delta_L$ at $\xi=0$ vanishes. By (b), $\Delta_L\le3\kappa+O(\xi^2)$ uniformly in $L$. Every lower bound $d_N$ of Corollary~\ref{cor:ym-orderN} is $3-64\xi+O(\xi^2)$, so its first-order loss is an artefact of the method.
\end{enumerate}
\end{proposition}
\begin{proof}
(a) The state $r_p$ is gauge invariant and orthogonal to $\psi$. By the min--max principle $\Delta_L\le q_A(r_p)/\|r_p\|^2$. By (F3), $q_A(r_p)=\kappa\langle\Gamma(W_p,W_p)\rangle_\rho=\kappa(3-\langle\chi_1\rangle_\rho)$ (Proposition~\ref{prop:kinetic}(a)). And $\|r_p\|^2=\Var_\rho(W_p)=1+\langle\chi_1\rangle_\rho-\langle W_p\rangle_\rho^2$, since $W_p^2=1+\chi_1(U_p)$.
(b) The two quantities in (a) are $(K_o)_{pp}\le3+w_K$ and $(G_0)_{pp}\ge1-w_0$.
(c) At $\xi=0$ the eigenvalue $3\kappa$ of $A$ on the physical space has the eigenspace spanned by the $W_p$: by Lemma~\ref{lem:casimir}, $c_j=3$ only for spin $\frac12$ on a cycle of length $4$, that is, on a plaquette. The first-order perturbation matrix is $\langle W_q,SW_{q'}\rangle=\sum_r\int W_qW_rW_{q'}\,dU=0$. Indeed, two distinct faces share at most one link, so in each such product some link occurs exactly once, in spin $\frac12$; and $\int W_q^3=0$. The ground energy has no first-order term either, since $\int S=0$. Finally, $w_K$ and $w_0$ are $O(\xi^2)$, and $d_N'(0)=-\frac32\ell_N'(0)+\frac32m_1-6t_1=-64$.
\end{proof}
\status{Proved here; (a) is unconditional, (b) rests on the inputs of Proposition~\ref{prop:bench}. The windows are computed by \note{ymr\_referee\_checks.py}.}

\subsection{Heat horizons, support quotients and the spatial limit}

\begin{itemize}
\item \emph{Finite heat time} (H33; the Gram comparison is the workbench's, the choice of $T$ is checked here). Replacing $\Phi$ by $\Phi_T=(I-e^{-TA})\Phi$ changes the Grams by factors between $(1-\varepsilon)^2$ and $1$, with $\varepsilon=e^{-\kappa d_{\mathrm{ph}}T}$. For a signed sum of $\log\det$ terms of total rank $B$ the error is at most $2B(-\log(1-\varepsilon))\le2B\varepsilon/(1-\varepsilon)$, and the choice $T=(\kappa d_{\mathrm{ph}})^{-1}\log(1+2B/\eta)$ makes $\varepsilon/(1-\varepsilon)=\eta/(2B)$, so the error is at most $\eta$.
\item \emph{Support quotients} (H31--H32; located). An exact isomorphism between quotients of inverse-energy spans and kernels, and the attained minimum of the family energy over changes of support.
\item \emph{The spatial limit at fixed spacing} (H34; located, an outline). Heat coefficients through degree $N$ agree on two boxes containing the same $N$-neighbourhood of the marks, so Theorem~\ref{thm:heat} makes the correlations converge as $L\to\infty$; the workbench then builds the infinite-volume semigroup, its invariant state and the relaxation bounds. This is a limit at fixed $a$ and $g$, not a continuum limit, as the workbench says.
\end{itemize}
